% Em Pelas Ondas do Rádio: Cultura Popular, Camponeses e o MEB analisa a participação de camponeses do nordeste brasileiro no Movimento de Educação de Base. A perspectiva da tese é a de demonstrar como os trabalhadores envolvidos com as escolas radiofônicas elaboraram ações para manutenção e reprodução da escola em sua comunidade, visando obter os benefícios necessários à reprodução e melhoria de seu modo de vida. A partir de representações políticas e culturais singulares, dentre as quais vigoraram: um sentido para escola, um papel para o sindicato e para participação política, preceitos do direito de uso da terra e dos direitos do trabalho, assim como, sentidos múltiplos para o uso do rádio como meio de comunicação, informação e lazer, os camponeses do MEB, foram coadjuvantes da proposição católica modernizadora de inícios de 1960. Isto posto, demarca que a ação do camponês nordestino e seu engajamento político, seja no MEB, nos sindicatos rurais, nas Juventudes Agrárias Católicas (JAC’s), no MCP, e nas mais diversas instâncias dos movimentos sociais do período, não se apartaram do processo modernizador. Neste sentido, considera-se que a modernização brasileira foi pauta das instituições, organismos políticos e partidos, assim como, do movimento social, instância em que ela foi ressignificada a partir de elementos da vida material, que envolviam diretamente, no momento em questão, a problemática do direito a terra, do direito a educação e cultura e dos direitos do trabalho.

Trata-se de dados ECG holter com duração de 24 horas de 3 derivações, coletados em pacientes com suspeita de arritmia cardíaca. O objetivo deste trabalho é desenvolver um algoritmo para detecção automática de eventos dado duração que o paciente esteve em observação médica, utilizando técnicas de aprendizado de máquina. Direfente de algumas pesquisa não irei utilizar técnica de classificação ou encontro de anomalia e nem detecção de series temporais, mas sim uma abordagem baseada em probabilidades de riscos para detecção de eventos. em vez de dizer se o paciente tem ou não arritmia, o objetivo é dizer qual a probabilidade do paciente ter um evento de uma parada cardiaca em um determinado período de tempo. com isso vou utilizar um método de Análise de Sobrevivência em Tempo Discreto (Logistic-Hazard) similar ao regressão de cox a diferença que são métodos diferentes porque a regressão de cox usa tempos contínuos e esse trabalho vai tratar de tempos discretos para análise de sobrevivência para prever o risco de eventos cardíacos em pacientes com suspeita de arritmia, com base em dados ECG holter de 24 horas. os resultados esperados são um modelo preditivo que possa auxiliar médicos na tomada de decisão clínica, identificando pacientes com maior risco de eventos cardíacos adversos. Os sinais de ECG vão entrar como covariáveis no modelo Logistic-Hazard(Risco Logistico) e o tempo até o evento (parada cardíaca) será a variável dependente. os ECGs vão passar pelo um pré-modelo de extração de características para extrair features relevantes dos sinais antes de serem inseridos no modelo de Análise de Sobrevivência em Tempo Discreto (Logistic-Hazard), como se trata de 3 derivações irei usar um modelo convolucional resnet de 18 camadas profundas por convulução 1D(1x3) para extrair recursos locais ao longo do tempo, junto com alguns dados de apoio para melhorar o aprendizado como idade, gênero e outras features. Como se trata de 3 sinais de derivações, vou utilizar uma rede neural convolucional de 1D para aprender sobre padrões dos sinais como intervalos de ondas como amplitude, anomalias, intervalos que será importante para o modelo Logistic-Hazard. 

% Separe as palavras-chave por ponto
\palavraschave{Logistic-Hazard. Resnet-18 1D. Educação de adultos. Escola rural.}